\lysh{21450}{2}
\begin{alist}
\item Ισχύει $x^2 + 4 > 0$ για κάθε πραγματικό αριθμό $x$, οπότε το πεδίο ορισμού της συνάρτησης $f$ είναι το $A_f = \mathbb{R}$. Η συνάρτηση $g$ ορίζεται για τους πραγματικούς αριθμούς $x$ για τους οποίους ισχύει $x > 0$. Άρα το πεδίο ορισμού της συνάρτησης $g$ είναι το $A_g = (0, +\infty)$.
\item Γνωρίζουμε ότι για $\alpha > 0$, $\alpha \neq 1$ και $x_1, x_2 > 0$ ισχύει η ισοδυναμία: $\log_\alpha x_1 = \log_\alpha x_2 \Leftrightarrow x_1 = x_2$. \\
Οπότε έχουμε:
\begin{align*}
f(x) &= g(x) \Leftrightarrow \\
\ln(x^2 + 4) &= \ln x + \ln 4 \Leftrightarrow \\
\ln(x^2 + 4) &= \ln(4x) \Leftrightarrow \\
x^2 + 4 &= 4x \Leftrightarrow \\
x^2 - 4x + 4 &= 0 \Leftrightarrow \\
(x - 2)^2 &= 0 \Leftrightarrow \\
x &= 2 > 0 \text{, δεκτή.}
\end{align*}
\end{alist}

